% ---------------------------------------------------------------------
% 附录 A：DQ/HDQ 层的验证性推导
% 附录 C：控制层的补充推导
% （原文档即无附录 B 与 C.5，此处保持一致）
% ---------------------------------------------------------------------

\section*{附录 A：DQ/HDQ 层的验证性推导}

\subsection*{A.1 定理 1 的 (i)(ii)(iii)}

(ii)：由 \eqref{eq:4.3} 直接右乘 $\tilde x$。(i)：翻转不变性：$2(-\dot{\tilde x})(-\tilde x^*)=2\dot{\tilde x}\tilde x^*$。(iii)：$\dot{\tilde x}=\dot{\hat{\underline x}}\hat{\underline x}_d^{\,*}+\hat{\underline x}\dot{\hat{\underline x}}_d^{\,*}$；用 $\dot{\hat{\underline x}}=\tfrac12\boldsymbol\xi\hat{\underline x}$ 与 $\dot{\hat{\underline x}}_d^{\,*}=-\tfrac12\hat{\underline x}_d^{\,*}\boldsymbol\xi_d$（后者由 $\dot{\hat{\underline x}}_d=\tfrac12\boldsymbol\xi_d\hat{\underline x}_d$ 取共轭）：
$\dot{\tilde x}=\tfrac12\boldsymbol\xi\tilde x-\tfrac12\hat{\underline x}\hat{\underline x}_d^{\,*}\boldsymbol\xi_d=\tfrac12\boldsymbol\xi\tilde x-\tfrac12\tilde x\boldsymbol\xi_d$。右乘 $2\tilde x^*$：$\tilde{\boldsymbol\xi}=\boldsymbol\xi-\tilde x\boldsymbol\xi_d\tilde x^*=\boldsymbol\xi-\mathrm{Ad}_{\tilde x}\boldsymbol\xi_d$。含速度级扰动 $\boldsymbol v_w,\boldsymbol v_c$（\cite{P2} 模型）时右端加 $\boldsymbol v_w+\boldsymbol v_c$。

\subsection*{A.2 引理 1 的证明}

\begin{equation*}
\tfrac{d}{dt}(\tilde x\boldsymbol a\tilde x^*)
=\dot{\tilde x}\boldsymbol a\tilde x^*+\tilde x\dot{\boldsymbol a}\tilde x^*+\tilde x\boldsymbol a\dot{\tilde x}^* .
\end{equation*}
代入 $\dot{\tilde x}=\tfrac12\tilde{\boldsymbol\xi}\tilde x$、$\dot{\tilde x}^*=-\tfrac12\tilde x^*\tilde{\boldsymbol\xi}$：
\begin{equation*}
=\tfrac12\tilde{\boldsymbol\xi}(\mathrm{Ad}_{\tilde x}\boldsymbol a)+\mathrm{Ad}_{\tilde x}\dot{\boldsymbol a}-\tfrac12(\mathrm{Ad}_{\tilde x}\boldsymbol a)\tilde{\boldsymbol\xi}
=\mathrm{Ad}_{\tilde x}\dot{\boldsymbol a}+\mathrm{ad}_{\tilde{\boldsymbol\xi}}(\mathrm{Ad}_{\tilde x}\boldsymbol a).
\end{equation*}
\eqref{eq:5.4}：对 \eqref{eq:4.4} 逐项求导并取 $\boldsymbol a=\boldsymbol\xi_d$。$\blacksquare$

\section*{附录 C：控制层的补充推导}

\subsection*{C.1 扰动项的适定性（\eqref{eq:5.1} 为何是显式等式，以及 $\alpha$ 条件的真正作用）}

$\boldsymbol\tau=\hat M\ddot{\boldsymbol q}_{\mathrm{ref}}+\hat C\dot{\boldsymbol q}+\hat g$ 代入真实动力学 $M\ddot{\boldsymbol q}+C\dot{\boldsymbol q}+\boldsymbol g=\boldsymbol\tau+\delta\boldsymbol\tau+\boldsymbol\tau_{\mathrm{ext}}$ 并解出 $\ddot{\boldsymbol q}$：
\begin{equation*}
\ddot{\boldsymbol q}=M^{-1}\hat M\ddot{\boldsymbol q}_{\mathrm{ref}}+M^{-1}\bigl(\Delta C\dot{\boldsymbol q}+\Delta\boldsymbol g+\delta\boldsymbol\tau+\boldsymbol\tau_{\mathrm{ext}}\bigr),
\qquad M^{-1}\hat M=I+M^{-1}\Delta M ,
\end{equation*}
即 \eqref{eq:5.1}。注意 $\ddot{\boldsymbol q}_{\mathrm{ref}}$ 由 \eqref{eq:5.2} 完全由\textbf{当前可测量} $(\boldsymbol q,\dot{\boldsymbol q},\hat{\underline x}_d,\boldsymbol\xi_d,\dot{\boldsymbol\xi}_d)$ 给出，\textbf{不依赖 $\ddot{\boldsymbol q}$}，故 \eqref{eq:5.1} 对 $\ddot{\boldsymbol q}$ 是\textbf{显式赋值}：不存在隐式代数环，既不需要移项、也不需要 Neumann 级数，更不会产生 $\alpha/(1-\alpha)$ 型的等效扰动增益因子。

由此可见 $\alpha<1/\mathrm{cond}_2(K_d)$（\eqref{eq:5.1f}）与适定性无关，而是为了使定理 3(d) 中的\textbf{有效阻尼保持正定}：乘性分量 $\Theta u_{\mathrm{fb}}$ 中与 $-K_de_\xi$ 同向的部分会削弱阻尼，由 $|e_\xi^\top\Theta(-K_de_\xi)|\le\alpha\lambda_{\max}(K_d)\|e_\xi\|^2$ 得 $\lambda_{\mathrm{eff}}=\lambda_{\min}(K_d)-\alpha\lambda_{\max}(K_d)$，正性恰好等价于 \eqref{eq:5.1f}。$\alpha$ 条件是一个小增益型的证书条件，而非微分方程适定性条件；若 $\alpha$ 过大，闭环仍然定义良好，但本文的 Lyapunov 证书失效。

\subsection*{C.2 定理 3(b) 的定量奇异值下界与收敛论证细节}

$\dot V=-e_\xi^\top K_de_\xi$ 只是关于全状态的\textbf{半}负定（$e_\xi=0$ 面上 $\dot V=0$），故不能直接给出收敛，必须补 LaSalle 步骤（定理 3(b) 证明第 4 步），而 LaSalle 步骤需要 $A$ 在工作域内可逆。定量化如下。

由 \eqref{eq:4.5}，$A=\begin{bmatrix}A_{11}&0\\ -[\mathcal T]_\times& I_3\end{bmatrix}$，$A_{11}=-\tfrac12(\tilde\eta I+[\mathcal O]_\times)$。其逆为
\begin{equation*}
\begin{gathered}
A^{-1}=\begin{bmatrix}A_{11}^{-1}&0\\ [\mathcal T]_\times A_{11}^{-1}& I_3\end{bmatrix},\\[2pt]
A_{11}^{-1}=-2\,\frac{\tilde\eta^2I+\tilde\eta[\mathcal O]_\times^{\top}+\mathcal O\mathcal O^\top}{\tilde\eta(\tilde\eta^2+\|\mathcal O\|^2)}
=-\frac2{\tilde\eta}\bigl(\tilde\eta^2I-\tilde\eta[\mathcal O]_\times+\mathcal O\mathcal O^\top\bigr),
\end{gathered}
\end{equation*}
用到 $\tilde\eta^2+\|\mathcal O\|^2=1$。$\|A_{11}^{-1}\|_2$ 无需放缩即可精确算出：记 $B\triangleq\tilde\eta I+[\mathcal O]_\times$（故 $A_{11}=-\tfrac12B$），则
\begin{equation*}
B^\top B=(\tilde\eta I-[\mathcal O]_\times)(\tilde\eta I+[\mathcal O]_\times)=\tilde\eta^2I-[\mathcal O]_\times^2=(\tilde\eta^2+\|\mathcal O\|^2)I-\mathcal O\mathcal O^\top=I-\mathcal O\mathcal O^\top ,
\end{equation*}
其特征值为 $\{\,\tilde\eta^2,\,1,\,1\,\}$（沿 $\mathcal O$ 方向为 $1-\|\mathcal O\|^2=\tilde\eta^2$，两个正交方向为 1）。于是 $B$ 的奇异值为 $\{\tilde\eta,1,1\}$，同时给出两个\textbf{精确}等式
\begin{equation*}
\|A_{11}\|_2=\tfrac12\sigma_{\max}(B)=\tfrac12 ,
\quad
\|A_{11}^{-1}\|_2=\frac2{\sigma_{\min}(B)}=\frac2{\tilde\eta} ,
\end{equation*}
前者正是 §5.2 与定理 3(b) 反复使用的精确值。再由分块下三角结构 $\|A^{-1}\|_2\le(1+\|\mathcal T\|)\allowbreak\|A_{11}^{-1}\|_2+1$，
\begin{equation*}
\sigma_{\min}(A)=\frac1{\|A^{-1}\|_2}\;\ge\;\Bigl[\frac{2(1+\|\mathcal T\|)}{\tilde\eta}+1\Bigr]^{-1} .
\end{equation*}
在 $\Omega_c$ 内 $\tilde\eta\ge\eta_0$（\eqref{eq:5.5c}）且 $\|\mathcal T\|\le\sqrt{2c/\lambda_{\min}(K_{p,T})}$，故 $\sigma_{\min}(A)\ge c_A(\eta_0,c)>0$ 一致成立。于是 LaSalle 第 4 步中由 $A^\top K_pe_z\equiv0$ 推出 $\|e_z\|\le\|K_p^{-1}\|_2\|A^{-\top}\|_2\cdot0=0$ 严格成立。局部指数率由线化矩阵 $F$ 的谱给出，逐分量具体极点见 \eqref{eq:5.8}。域限制 $\tilde\eta>0$ 与 \cite{P2} Remark 1 的 unwinding 条件一致。

\begin{remark}
注：无扰情形下 LaSalle 路线自身是完整的，但不能改用级联 ISS 论证闭合。
\end{remark}

\subsection*{C.3 定理 3(c) 的二次型/Schur 补细节}

\textbf{(c-1) Schur 补判据。} 记供给率 $s(e_\xi,d)\triangleq-\tfrac1{2\kappa}\|e_\xi\|^2+\tfrac{\gamma_a^2}2\|d\|^2$。性能目标 $\dot V\le s$ 等价于 $-e_\xi^\top K_de_\xi+e_\xi^\top d-s(e_\xi,d)\le0$，整理为 $-[e_\xi;d]^\top M[e_\xi;d]\le0$，其中 $M$ 为 \eqref{eq:5.6a} 的分块矩阵。$M\succeq0$ 且右下块 $\tfrac{\gamma_a^2}2I\succ0$，取 Schur 补得
\begin{equation*}
M\succeq0\iff K_d-\tfrac1{2\kappa}I-\tfrac1{2\gamma_a^2}I\succeq0\iff K_d\succeq\tfrac12(\kappa^{-1}+\gamma_a^{-2})I .
\end{equation*}
这是\textbf{当且仅当}判据：不定号交叉项 $e_\xi^\top d$ 保留在二次型内整体判定，无任何符号放缩。

\textbf{(c-2) 两处恒零的几何含义与失效条件。} 拆分精确成立依赖两条叉积混合积恒等式：(i) $(A^\top e_z)_\omega$ 中 $[\mathcal T]_\times\mathcal T=\mathcal T\times\mathcal T=0$；(ii) $\dot{\mathcal T}=-[\mathcal T]_\times\tilde\omega+\tilde v$ 中耦合项对 $\tfrac12\|\mathcal T\|^2$ 不做功，$\mathcal T^\top(\mathcal T\times\tilde\omega)=0$。二者是代数恒等式，不依赖工作域，故 (c-2) 与 (c-1) 一样是全局结论。若 $K_d$ 非块对角，$-e_\xi^\top K_de_\xi$ 含 $\tilde\omega^\top K_{\omega v}\tilde v$ 型交叉项，两分量能量不再分离，退回合并判据 (c-1)。

\textbf{$K_{p,T}$ 各向同性的必要性。} 将平移刚度写作一般对称正定 $K_{p,T}$ 时，上述两处恒零均失效：(i) $(A^\top K_pe_z)_\omega=A_{11}^\top K_{p,O}\mathcal O+[\mathcal T]_\times K_{p,T}\mathcal T$，而 $[\mathcal T]_\times K_{p,T}\mathcal T=\mathcal T\times(K_{p,T}\mathcal T)\ne0$ 除非 $\mathcal T$ 是 $K_{p,T}$ 的特征向量；(ii) 平移储能变为 $\tfrac12\mathcal T^\top K_{p,T}\mathcal T$，其导数中的耦合项 $-\mathcal T^\top K_{p,T}[\mathcal T]_\times\tilde\omega$ 不再为零（仅当 $K_{p,T}=k_{p,T}I_3$ 时由 $\mathcal T^\top[\mathcal T]_\times=0$ 而消失）。两项残留均是 $\tilde\omega$--$\mathcal T$ 型耦合，使 $V_\omega,V_v$ 不再各自闭合，退回 (c-1)。

\subsection*{C.4 strictification 与局部 ISS（未来工作）}

定理 3(d) 的三条缺口的共同根源是 $\dot V$ 中没有 $-\|e_z\|^2$ 型负项。标准补救是 strictification：取 $W=V+\epsilon e_z^\top K_pA(\tilde x)e_\xi$（$\epsilon>0$ 待定），在 $\Omega_c$ 内 $\epsilon$ 足够小时 $W$ 与 $V$ 等价，而求导后新增 $-\epsilon e_z^\top K_pAA^\top K_pe_z$ 项提供 $e_z$ 方向负定。若能验证剩余交叉项（含 $\dot A$ 项）可被两个负定项吸收，则 $\dot W\le-c_1\|(e_z,e_\xi)\|^2+c_2\|d\|\cdot\|(e_z,e_\xi)\|$，得到真正的局部 ISS-Lyapunov 函数。本文未完成此验证（关键难点是 $\dot A$ 项的一致界与 $\epsilon$ 的可行区间非空性），列为 §7 局限 (i)。

\subsection*{C.6 定理 3(b) 的完整证明}

\begin{enumerate}[leftmargin=2em]
  \item \textbf{交叉项精确相消}：沿 \eqref{eq:5.5}（$d\equiv0$），
  \begin{equation*}
  \dot V=e_\xi^\top\dot e_\xi+e_z^\top K_p\dot e_z
  =e_\xi^\top\bigl(-K_de_\xi-A^\top K_pe_z\bigr)+e_z^\top K_pAe_\xi
  =-e_\xi^\top K_de_\xi ,
  \end{equation*}
  因 $e_z^\top K_pAe_\xi=(A^\top K_p^\top e_z)^\top e_\xi=(A^\top K_pe_z)^\top e_\xi$——这里\textbf{只}用到 $K_p$ 对称与 $K_p$ 写在 $A^\top$ 内侧，不需要 $K_p$ 为标量。故 $\dot V\le0$，$\Omega_c$ 正向不变；又 $V\le c$ 给出 $\|e_\xi\|\le\sqrt{2c}$、$\|e_z\|\le\sqrt{2c/\lambda_{\min}(K_p)}$，故 $\Omega_c$ 紧。

  \item \textbf{工作域保号 \eqref{eq:5.5c}}：$V\le c$ 蕴含 $\tfrac12\mathcal O^\top K_{p,O}\mathcal O\le c$，即 $\|\mathcal O\|^2\le2c/\lambda_{\min}(K_{p,O})<1$（由 $c<c^*$）。由 $\mathcal O=-\mathrm{Im}\,\tilde r$ 与 $\|\tilde r\|=1$ 得 $\tilde\eta^2+\|\mathcal O\|^2=1$，故 $|\tilde\eta|\ge\eta_0>0$；$\tilde\eta(t)$ 连续且恒不为零，符号不可突变，由 $\tilde\eta(0)>0$ 得 \eqref{eq:5.5c}。

  \item \textbf{$A$ 的行列式}：$A$ 块下三角（\eqref{eq:4.5}），故
  \begin{equation*}
  \det A=\det\bigl(-\tfrac12(\tilde\eta I_3+[\mathcal O]_\times)\bigr)\cdot\det I_3
  =\bigl(-\tfrac12\bigr)^3\tilde\eta\bigl(\tilde\eta^2+\|\mathcal O\|^2\bigr)=-\tfrac18\tilde\eta ,
  \end{equation*}
  用到 $\det(aI_3+[b]_\times)=a(a^2+\|b\|^2)$ 与 $\tilde\eta^2+\|\mathcal O\|^2=1$。定量奇异值下界 $\sigma_{\min}(A)\ge\bigl[2(1+\|\mathcal T\|)/\tilde\eta+1\bigr]^{-1}$ 见附录 C.2。

  \item \textbf{LaSalle}：$\Omega_c$ 紧且不变，$E\triangleq\{\dot V=0\}\cap\Omega_c=\{e_\xi=0\}\cap\Omega_c$。若轨迹全程留在 $E$ 内：$e_\xi\equiv0\Rightarrow\dot e_\xi\equiv0\Rightarrow A^\top K_pe_z\equiv0$；由第 3 步 $A$ 可逆、$K_p\succ0$ 得 $e_z\equiv0$。故 $E$ 内最大不变集为 $\{(0,0)\}$，由 LaSalle 不变集定理（\cite{Kha02} Thm 4.4）得 (iii)。

  \item \textbf{局部指数稳定}：在 $(0,0)$ 处 $\tilde x\to1$，$A\to A_0=\mathrm{diag}(-\tfrac12I_3,I_3)$，\eqref{eq:5.5} 的雅可比为
  \begin{equation*}
  F=\begin{bmatrix}0_6 & A_0\\ -A_0^\top K_p & -K_d\end{bmatrix}.
  \end{equation*}
  对该 LTI 系统同一个 $V$ 仍给出 $\dot V=-e_\xi^\top K_de_\xi\le0$，且 $A_0$ 可逆，重复第 4 步得 LTI 系统渐近稳定，故 $F$ 为 Hurwitz；再由 Lyapunov 线化定理（\cite{Kha02} Thm 4.7）得非线性系统在原点邻域指数稳定。块对角 $K_d,K_p$ 下 $F$ 逐分量解耦，其两个二阶多项式与极点由 \eqref{eq:5.8} 显式给出。
\end{enumerate}

\subsection*{C.7 定理 3(c) 的完整证明}

\textbf{第一步（$\dot V$ 精确式）}：含扰时定理 3(b) 证明第 1 步的交叉项相消与 $d$ 无关，故
\begin{equation*}
\dot V=-e_\xi^\top K_de_\xi+e_\xi^\top d
\end{equation*}
精确成立（此式不含任何放缩；$e_\xi^\top d$ 可正可负）。

\textbf{第二步（(c-1) 判据的等价性）}：性能目标即 $-e_\xi^\top K_de_\xi+e_\xi^\top d+\tfrac1{2\kappa}\|e_\xi\|^2-\tfrac{\gamma_a^2}2\|d\|^2\le0$，等价于二次型不等式
\begin{equation*}
\begin{bmatrix}e_\xi\\ d\end{bmatrix}^{\!\top}\!M\!
\begin{bmatrix}e_\xi\\ d\end{bmatrix}\ge0\quad\forall(e_\xi,d)\in\mathbb R^{12},
\end{equation*}
即 $M\succeq0$。右下块 $\tfrac{\gamma_a^2}2I\succ0$，取 Schur 补得 $K_d-\tfrac1{2\kappa}I-\tfrac1{2\gamma_a^2}I\succeq0$，即 \eqref{eq:5.6a}。关键在于不定号交叉项 $e_\xi^\top d$ 保留在二次型内整体判定，不经任何符号放缩。

\textbf{第三步（全局存在性）}：由 $M\succeq0$ 得 $\dot V\le\tfrac{\gamma_a^2}2\|d\|^2$，故 $V(t)\le V(0)+\tfrac{\gamma_a^2}2\|d_{L_2}\|_{L_2}^2<\infty$，$(e_z,e_\xi)$ 一致有界，解在 $[0,\infty)$ 上存在（无有限时间逃逸）。

\textbf{第四步（积分收尾）}：在 $[0,T]$ 上积分 $\dot V\le-\tfrac1{2\kappa}\|e_\xi\|^2+\tfrac{\gamma_a^2}2\|d\|^2$，弃去 $V(T)\ge0$，令 $T\to\infty$（单调收敛）即得 \eqref{eq:5.6}。

\textbf{第五步（(c-2) 分量解耦）}：块对角 $K_d$ 与各向同性平移刚度 $K_{p,T}=k_{p,T}I_3$ 下两分量储能精确解耦，关键是两处混合积恒零：由 \eqref{eq:4.5}，$(A^\top K_pe_z)_\omega=A_{11}^\top K_{p,O}\mathcal O+k_{p,T}[\mathcal T]_\times\mathcal T=A_{11}^\top K_{p,O}\mathcal O$（$\mathcal T\times\mathcal T=0$，故旋转反馈不含 $\mathcal T$）、$(A^\top K_pe_z)_v=k_{p,T}\mathcal T$；又 $\dot{\mathcal T}=-[\mathcal T]_\times\tilde\omega+\tilde v$ 中的耦合项做功为零（$\mathcal T\cdot(\mathcal T\times\tilde\omega)=0$）。于是位姿交叉项在两分量内分别由 $K_{p,O}$ 对称与 $k_{p,T}$ 为标量而精确相消，
\begin{equation*}
\dot V_\omega=-\tilde\omega^\top K_\omega\tilde\omega+\tilde\omega^\top d_\omega,
\qquad
\dot V_v=-\tilde v^\top K_v\tilde v+\tilde v^\top d_v ,
\end{equation*}
对每个分量重复第二至第四步的论证（$I_6\to I_3$）即得 \eqref{eq:5.6b}$\Rightarrow$\eqref{eq:5.6p}。两处恒零为代数恒等式，故 (c-1)/(c-2) 的全部结论均不依赖工作域 $\tilde\eta>0$。

\subsection*{C.8 定理 3(d) 的完整证明}

\begin{enumerate}[leftmargin=2em]
  \item \textbf{精确耗散等式与乘法项展开}：由定理 3(c) 证明第一步，$\dot V=-e_\xi^\top K_de_\xi+e_\xi^\top d$ 精确成立（无任何放缩）。代入 \eqref{eq:5.1d} 与 \eqref{eq:5.2} 的 $u_{\mathrm{fb}}=-K_de_\xi-A^\top K_pe_z$：
  \begin{equation*}
  \dot V=-e_\xi^\top K_de_\xi\;\underbrace{-\,e_\xi^\top\Theta K_de_\xi}_{\text{与阻尼同类}}\;\underbrace{-\,e_\xi^\top\Theta A^\top K_pe_z}_{\text{与扰动同类}}\;+\;e_\xi^\top d_{\mathrm{ex}} .
  \end{equation*}

  \item \textbf{两类乘法项的分别回收}：$|e_\xi^\top\Theta K_de_\xi|\le\alpha\lambda_{\max}(K_d)\|e_\xi\|^2$（回收进阻尼，得 \eqref{eq:5.7a} 的 $\lambda_{\mathrm{eff}}$；由 (A3)/\eqref{eq:5.1f} 即 $\alpha<\lambda_{\min}(K_d)/\lambda_{\max}(K_d)$ 得 $\lambda_{\mathrm{eff}}>0$）；$|e_\xi^\top\Theta A^\top K_pe_z|\allowbreak\le\alpha\|A\|_2\allowbreak\lambda_{\max}(K_p)\|e_z\|\,\|e_\xi\|$（回收进等效扰动幅值 $D$）。

  \item \textbf{$A$ 的谱范数界}：由 $A_{11}^\top A_{11}=\tfrac14(I_3-\mathcal O\mathcal O^\top)$ 得 $\|A_{11}\|_2=\tfrac12$（\textbf{精确值}，与 $\tilde x$ 无关；奇异值计算见附录 C.2），再由 $A$ 的块下三角结构得 $\|A\|_2\le\max\{\|A_{11}\|_2,1\}+\|[\mathcal T]_\times\|_2=1+\|\mathcal T\|$。在 $\Omega_c$ 上 $\|\mathcal T\|\le\sqrt{2c/\lambda_{\min}(K_{p,T})}$、$\|e_z\|\le\sqrt{2c/\lambda_{\min}(K_p)}$，代入第 2 步即得 \eqref{eq:5.7b} 与
  \begin{equation}
  \dot V\ \le\ -\lambda_{\mathrm{eff}}\|e_\xi\|^2+D\,\|e_\xi\| .
  \tag{5.7d}\label{eq:5.7d}
  \end{equation}

  \item \textbf{Young 与积分收尾}：$D\|e_\xi\|\le\tfrac{\lambda_{\mathrm{eff}}}2\|e_\xi\|^2+\tfrac{D^2}{2\lambda_{\mathrm{eff}}}$，故 $\dot V\le-\tfrac{\lambda_{\mathrm{eff}}}2\|e_\xi\|^2+\tfrac{D^2}{2\lambda_{\mathrm{eff}}}$。在 $[0,T]$ 上积分并弃去 $V(T)\ge0$：
  \begin{equation*}
  \begin{gathered}
  \tfrac{\lambda_{\mathrm{eff}}}2\int_0^T\|e_\xi\|^2dt\\[2pt]
  \ \le\ V(0)+\tfrac{D^2}{2\lambda_{\mathrm{eff}}}\,T ,
  \end{gathered}
  \end{equation*}
  两端除以 $\tfrac{\lambda_{\mathrm{eff}}}2T$ 即 \eqref{eq:5.7c}；令 $T\to\infty$ 得 \eqref{eq:5.7}。$\alpha\to0$ 时 $D\to D_{\mathrm{ex}}$、$\lambda_{\mathrm{eff}}\to\lambda_{\min}(K_d)$，退化为经典形式。
\end{enumerate}
